\documentclass[../../main.tex]{subfiles}
\begin{document}

\begin{center}
  {\fontsize{48pt}{58pt}\selectfont \textbf{1965}\normalsize\textbf{年度}}
\end{center}

\begin{table}[htb]
  \centering
  \caption{出題テーマと難易度}
  \begin{tabular}{|c|p{9cm}|c|}
    \hline
    \textbf{問題番号} & \textbf{テーマ} & \textbf{難易度} \\
    \hline
    第1問 & 三種類の新聞の購読率と集合 & 易 \\
    \hline
    第2問 & 山頂の方位と仰角から求める高さ & 普 \\
    \hline
    第3問 & 双曲線の接線と囲まれる面積 & 易 \\
    \hline
    第4問 & 半円内の四辺形の面積の最大化 & 普 \\
    \hline
    第5問 & 三角関数の曲線に引いた平行な接線 & 普 \\
    \hline
    第6問 & 円弧で囲まれる領域の回転体 & 易 \\
    \hline
  \end{tabular}
\end{table}

\begin{table}[htb]
  \centering
  \caption{解答}
  \renewcommand{\arraystretch}{1.4}
  \begin{tabular}{|c|c|p{9cm}|}
    \hline
    \textbf{問題番号} & \textbf{小問} & \textbf{答え} \\
    \hline
    第1問 & (1) & $48\%$ \\
    \cline{2-3}
    & (2) & $10\%$ \\
    \cline{2-3}
    & (3) & $5\%$ \\
    \cline{2-3}
    & (4) & $2\%$ \\
    \hline
    第2問 & -- & 海抜$1954\,\mathrm{m}$ \\
    \hline
    第3問 & -- & 最大面積は$\dfrac{2}{3}$ \\
    \hline
    第4問 & -- & $0<k\le\sqrt{2}-1$なら$OP=\dfrac{k+1}{\sqrt{2}}$，$k\ge\sqrt{2}-1$なら$P=C$ \\
    \hline
    第5問 & -- & 接点の$x$座標は$\dfrac{\pi}{2},\dfrac{\pi}{3},\dfrac{2\pi}{3}$．各接線の方程式は解答参照 \\
    \hline
    第6問 & -- & 体積は$\dfrac{2\pi}{3}$ \\
    \hline
  \end{tabular}
\end{table}

\paragraph{試験形態}
6問．試験時間は要確認．

\end{document}
